% ============================================================
%  Sage: Academic Paper — v2 draft
%  Template: generic two-column conference style
%  To switch venue:
%    ACM:  \documentclass[sigconf]{acmart}
%    IEEE: \documentclass[conference]{IEEEtran}
% ============================================================
\documentclass[10pt,twocolumn]{article}

\usepackage[margin=0.75in,top=1in,bottom=1in]{geometry}
\usepackage{times}
\usepackage{graphicx}
\usepackage{amsmath}
\usepackage{booktabs}
\usepackage{tabularx}
\usepackage{array}
\usepackage{enumitem}
\usepackage{listings}
\usepackage{xcolor}
\usepackage{hyperref}
\usepackage{microtype}
\usepackage{caption}
\usepackage{abstract}
\usepackage{titlesec}
\usepackage{cite}

% Section formatting
\titleformat{\section}{\normalfont\large\bfseries}{\thesection.}{0.5em}{}
\titleformat{\subsection}{\normalfont\normalsize\bfseries}{\thesubsection}{0.5em}{}
\titlespacing*{\section}{0pt}{10pt}{4pt}
\titlespacing*{\subsection}{0pt}{8pt}{3pt}

% Code listing style
\lstset{
  basicstyle=\ttfamily\footnotesize,
  backgroundcolor=\color{gray!10},
  frame=single,
  framerule=0.4pt,
  rulecolor=\color{gray!50},
  breaklines=true,
  captionpos=b,
  numbers=none,
}

% Compact list spacing
\setlist[itemize]{noitemsep, topsep=3pt, leftmargin=*}
\setlist[enumerate]{noitemsep, topsep=3pt, leftmargin=*}

\hypersetup{colorlinks=true, linkcolor=blue!60!black, citecolor=blue!60!black, urlcolor=blue!60!black}

% ============================================================
\begin{document}

\title{\textbf{Sage: A Traceable, Shareable Natural Language\\
Agent for Scientific Data Exploration}}

\author{
[Author Names --- to be completed]\\
\small [Affiliation]\\
\small [Target Conference] $\bullet$ [Year]
}

\date{}
\maketitle

% ============================================================
\begin{abstract}
[To be written after the body sections are settled.]
\end{abstract}

% ============================================================
\section{Introduction}

AI coding agents---Claude Code, GitHub Copilot Workspace, Devin, and their kin---have
transformed software engineering. A developer states a task in plain English; the agent
writes, runs, and debugs code autonomously; the deliverable is working software. For
software engineering, this model is a genuine breakthrough.

Data scientists, however, do not need working software---they need \emph{answers}.
A researcher studying earthquake impacts does not want a Python script; they want a map
showing where the M5.7 event occurred, a table of aftershocks, and a chart of GPS
displacement around the event date. A hydrologist assessing a flood scenario does not
want code that reads a raster---they want a breakdown of flooded schools, commercial
buildings, and elderly residents at risk. The deliverable is scientific insight: visual,
structured, and interpretable. Coding agents stop at the script. Scientists need to go
further.

\textbf{Sage} (Science Agent) is a natural language data science agent that extends the
coding agent model to its logical conclusion for scientific work. Like a coding agent,
Sage takes a plain English prompt and autonomously generates and executes code behind the
scenes. Unlike a coding agent, Sage does not surface that code as the result. Instead,
each prompt triggers a complete end-to-end pipeline---data discovery, retrieval,
processing, and analysis---and delivers the outcome directly as interactive maps,
annotated charts, structured tables, and a natural language summary, all rendered inline
in a standard Jupyter notebook. The code is a hidden implementation detail. The answer is
the deliverable.

A useful mental model is \emph{ChatGPT for scientific data}: the user asks a question in
plain English and receives a rich, visual, data-backed answer---not a script to run
separately. But Sage goes beyond conversational AI in two critical ways. First, every
step the agent takes is logged in a collapsible inspection panel embedded directly in the
notebook cell that triggered it: the exact API call issued, the code generated, the files
produced, the intermediate results examined. Nothing is a black box. Second, the result
lives in a standard \texttt{.ipynb} file---shareable, archivable, and renderable by any
Jupyter environment. A colleague who receives a Sage notebook sees not just the final
answer but the complete, step-by-step analytical record that produced it.

Sage is deployed as a JupyterHub magic command (\texttt{\%\%ask}) with no per-notebook
setup required. It is LLM-agnostic---supporting OpenAI, Anthropic, Google, and any
OpenAI-compatible endpoint---and runs on any JupyterHub installation or local JupyterLab
environment. Domain knowledge is supplied through a plug-in skill architecture: plain-text
\texttt{SKILL.md} files that teach the agent about specific data sources, APIs, and
scientific workflows. Skills ship with the Docker image for pre-configured deployments,
and can also be discovered and installed live from the SkillsMP community
marketplace~\cite{skillsmp} within a running session, without restarting the kernel.

This paper makes the following contributions:

\begin{itemize}
  \item \textbf{A data science agent model} that extends coding agents beyond code
        generation to deliver full scientific results---maps, charts, tables, and
        narrative---with all intermediate steps transparently logged and inspectable.

  \item \textbf{Cross-cell conversational memory} maintained as a Python-native in-kernel
        message list, enabling coherent multi-step reasoning across notebook cells with
        no external database or infrastructure.

  \item \textbf{A two-tier skill architecture} combining image-bundled domain skills
        for pre-configured deployments with runtime skill installation from the SkillsMP
        marketplace, all operable in natural language within a live session.

  \item \textbf{Notebook-scoped persistent output management} that automatically
        dispatches agent-generated files (GeoJSON, WMS, CSV, PNG) into multi-layer
        interactive maps and inline visualizations embedded in the final report.

  \item \textbf{Three end-to-end scientific demonstrations} spanning seismic event
        analysis with GNSS geodetic data, multi-source flood impact assessment, and
        on-the-fly skill installation and immediate expert-level use.
\end{itemize}

The remainder of this paper is organized as follows. Section~2 surveys related work.
Section~3 describes the Sage architecture. Section~4 presents the three demonstrations.
Section~5 discusses design principles. Section~6 addresses limitations and future work.
Section~7 concludes.

% ============================================================
\section{Related Work}

[To be written.]

% ============================================================
\section{System Architecture}

[To be written.]

% ============================================================
\section{Demonstrations}

[To be written.]

% ============================================================
\section{Design Principles and Discussion}

[To be written.]

% ============================================================
\section{Limitations and Future Work}

[To be written.]

% ============================================================
\section{Conclusion}

[To be written.]

% ============================================================
\section*{Acknowledgments}

[To be completed.]

% ============================================================
\bibliographystyle{plain}
\begin{thebibliography}{99}

\bibitem{skillsmp}
SkillsMP. \textit{The AI Agent Skills Marketplace}. \url{https://skillsmp.com}

\bibitem{sage_github}
Sage GitHub Repository. \url{https://github.com/klinucsd/sage}

\end{thebibliography}

\end{document}
